\Title{Кибизов~К.О., Сальников~А.Н.}{Технические аспекты выбора пути передачи данных через физическую среду в~кластере с~графическими процессорами}

\section*{Введение}

В вычислительном кластере данные между графическими процессорами (GPU) могут передаваться через разные физические среды: одни соединяют устройства внутри узла, другие~--- узлы между собой. Поэтому между двумя GPU обычно существует несколько возможных путей передачи. Например, передача данных может осуществляться как напрямую между устройствами, так и с промежуточным копированием через оперативную память. Такие варианты передачи возможны как между GPU в пределах одного вычислительного узла, так и между GPU и сетевым адаптером при межузловой передаче.

Выбор пути зависит не только от аппаратуры, но и от программного стека: операционной системы, драйверов устройств, коммуникационных библиотек и протоколов передачи. Программа взаимодействует с этим стеком через соответствующие интерфейсы: CUDA или OpenCL используются для работы с GPU, а MPI~--- для передачи данных между процессами. Решения, от которых зависит путь, принимаются на разных уровнях стека, поэтому в сложной среде путь не очевиден даже при известной топологии. От выбранного пути зависит задержка передачи, то есть время, необходимое для доставки данных от одного GPU к другому.

Для разных пар компонентов среды передачи данных задержка может заметно различаться. Неоднородность задержек между парами узлов кластера показана в~[\Refn{salnikov2017}], а в~[\Refn{begaev2020}] предложен метод их измерения для GPU, расположенных на разных узлах. Настоящая работа продолжает этот подход и расширяет набор тестов проекта Clustbench в ветке \texttt{gpu\_tests}~[\Refn{clustbench}]. Ранее на каждом узле запускался один MPI-процесс, обслуживавший все его GPU: передача внутри узла измерялась средствами CUDA без участия MPI, а между узлами всегда шла через оперативную память. Теперь каждому GPU соответствует отдельный MPI-процесс, поэтому любая передача, как внутри узла, так и между узлами, измеряется через MPI, а при необходимости передачу можно принудительно вести через оперативную память.

Цель работы~--- выяснить, как настройки Open MPI, транспортного слоя UCX и планировщика Slurm определяют путь передачи данных между GPU, и исследовать задержки передачи в зависимости от размера сообщения, взаимного расположения GPU и среды передачи.

\section{Модель коммуникационной среды и постановка задачи}

Коммуникационная среда представляется в виде графа. Его вершины соответствуют аппаратным компонентам: процессорам, GPU, сетевым адаптерам и коммутаторам. Рёбра соответствуют физическим соединениям: PCIe, NVLink, межпроцессорному интерфейсу, Ethernet, InfiniBand. При тестировании GPU передают друг другу сообщения~--- синтетически сформированными блоками данных заданного размера. При передаче сообщение проходит по графу через последовательность вершин и рёбер, то есть по некоторому пути. Множество возможных путей зависит от топологии и программного окружения, а конкретный путь выбирает транспортный слой. Пример такой среды~--- узел с несколькими GPU~--- показан на рис.~1.

\medskip
\noindent\begin{minipage}{\textwidth}
\includegraphics[width=\textwidth]{kibizov_421_arch}
\begin{center}
Рис.~1: Пример архитектуры узла с несколькими GPU (CPU~--- процессор, RAM~--- оперативная память, NIC~--- сетевой адаптер, switch~--- коммутатор)
\end{center}
\end{minipage}
\medskip

Основные пути передачи данных между GPU показаны на рис.~2. Без промежуточного копирования данные передаются внутри узла средствами GPUDirect Peer-to-Peer~[\Refn{cudap2p}] (путь a), между узлами~--- GPUDirect RDMA~[\Refn{gpudirect}] (путь c). В остальных случаях данные копируются через оперативную память (пути b и d).

Предметом изучения служит сквозная задержка передачи сообщения между упорядоченной парой различных устройств, то есть время от начала передачи на устройстве-отправителе до завершения приёма на устройстве-получателе.

\medskip
\noindent\begin{minipage}{\textwidth}
\includegraphics[width=\textwidth]{kibizov_421_paths}
\begin{center}
Рис.~2: Пути передачи данных между GPU (RAM~--- оперативная память, NIC~--- сетевой адаптер)
\end{center}
\end{minipage}
\medskip

\section{Настройка коммуникационного стека}

Путь передачи задаётся на нескольких уровнях: аппаратуры, драйверов, программного обеспечения и самой программы. Ниже настройки рассмотрены по этим уровням.

\subsection{Уровень аппаратуры}

На уровне аппаратуры путь определяется расположением устройств в узле: к какому процессору и коммутатору PCIe подключены GPU и сетевые адаптеры и соединены ли GPU через NVLink. Топологию узла показывает команда \texttt{nvidia-smi topo -m}.

При передаче без промежуточного копирования данные идут между GPU и сетевым адаптером мимо процессора, поэтому адаптер нужно выбирать по близости к GPU в топологии PCIe, а не к ядрам процесса, как это по умолчанию делает UCX.

\subsection{Уровень драйверов}

Прямой доступ GPU к памяти другого GPU того же узла (GPUDirect Peer-to-Peer~[\Refn{cudap2p}]) обеспечивает драйвер NVIDIA. Разрешён ли он для пары GPU, показывает команда \texttt{nvidia-smi topo -p2p}. Для передачи между узлами без промежуточного копирования (GPUDirect RDMA~[\Refn{gpudirect}]) сетевой адаптер должен обращаться к памяти GPU, для чего должен быть загружен модуль ядра \texttt{nvidia\_peermem}. На рассмотренных узлах оба условия выполнены.

\subsection{Уровень программного обеспечения}

\textbf{Архитектура UCX.} Open MPI~[\Refn{openmpi}] передаёт сообщения через транспортный слой UCX~[\Refn{ucx}]. UCX состоит из двух уровней. Нижний уровень предоставляет транспорты~--- способы передачи для конкретной аппаратуры (таблица~1). Верхний уровень реализует протоколы передачи сообщений и для каждой пары процессов выбирает, какими транспортами пользоваться.

\textbf{Параметры сборки.} UCX 1.15.0 собран с поддержкой CUDA, InfiniBand и CMA, но без библиотеки GDRCopy, ускоряющей копирование небольших сообщений из памяти GPU. Open MPI 5.0.10 собран с этой версией UCX и с поддержкой CUDA. Использовался CUDA Toolkit 12.8.

\textbf{Параметры Open MPI и UCX.} В обеих средах передачи, \texttt{auto} и \texttt{host} (раздел~3), Open MPI передаёт сообщения через UCX (\texttt{OMPI\_MCA\_pml=ucx}) с включённой поддержкой CUDA (\texttt{OMPI\_MCA\_opal\_cuda\_support=1}). Набор доступных транспортов задаётся переменной \texttt{UCX\_TLS}. В среде \texttt{auto} в него входят транспорты CUDA, а переменная \texttt{UCX\_IB\_GPU\_DIRECT\_RDMA=y} разрешает сетевому адаптеру обращаться к памяти GPU. В среде \texttt{host} в MPI передаются только указатели на оперативную память, поэтому транспорты CUDA исключены, а GPUDirect RDMA запрещена. Если нужный транспорт не указан или недоступен, UCX не сообщает об ошибке, а передаёт данные другим путём, как правило через оперативную память.

\medskip
\noindent\begin{minipage}{\textwidth}
\begin{center}
\small
\renewcommand{\arraystretch}{1.15}
\begin{tabular}{@{}l>{\raggedright\arraybackslash}p{0.5\textwidth}cc@{}}
\hline
Транспорт & Назначение & \texttt{auto} & \texttt{host} \\
\hline
\texttt{cuda\_copy} & копирование между памятью GPU и оперативной памятью & + & --- \\
\texttt{cuda\_ipc} & доступ к памяти GPU другого процесса узла (GPUDirect P2P) & + & --- \\
\texttt{rc}, \texttt{ud}, \texttt{dc} & передача между узлами по сети InfiniBand & + & + \\
\texttt{cma} & копирование между процессами узла средствами ядра ОС & + & + \\
\texttt{sm} & передача через разделяемую память узла & + & + \\
\texttt{self} & передача внутри процесса & + & + \\
\hline
\end{tabular}

\smallskip
{\normalsize Таблица~1: Транспорты UCX в средах \texttt{auto} и \texttt{host}}
\end{center}
\end{minipage}
\medskip

Сообщение передаётся одним из двух протоколов. В протоколе eager отправитель сразу посылает данные, не дожидаясь готовности получателя, и получатель при необходимости складывает их во временный буфер. Сообщение в памяти GPU при этом копируется через оперативную память (пути b и d на рис.~2). В протоколе rendezvous отправитель сначала посылает только запрос, и данные передаются после того, как получатель подготовит буфер приёма. Это требует дополнительного обмена служебными сообщениями, но позволяет передать данные без промежуточного копирования (пути a и c). Обычно протокол eager выгоден для небольших сообщений, поскольку не требует согласования, а rendezvous~--- для крупных. Однако для сообщений в памяти GPU eager требует копирования через оперативную память, и его преимущество может теряться (раздел~5.1).

Выбор протокола определяется порогом \texttt{UCX\_RNDV\_THRESH}, который задаётся при запуске: сообщения меньше порога передаются протоколом eager, остальные~--- протоколом rendezvous. Использовались три значения порога. Значение по умолчанию UCX выбирает автоматически. При значении 0 любое сообщение передаётся протоколом rendezvous. Третье значение выбрано больше наибольшего размера сообщения в измерениях, поэтому все сообщения передаются протоколом eager. Два крайних значения позволяют сравнить оба протокола на всём диапазоне размеров.

\textbf{Планировщик Slurm.} Процессы запускаются командой \texttt{srun} с интерфейсом PMIx, по одному на GPU. Флаг \texttt{-{}-gres-flags=enforce-binding} требует, чтобы ядра процесса находились на том же сокете, что и его GPU~[\Refn{slurm}], иначе передача данных между процессором и GPU пойдёт через межпроцессорный интерфейс. Ядра закрепляются за процессом параметром \texttt{-{}-cpu-bind=cores}. GPU запрашиваются на узел (\texttt{-{}-gpus-per-node}), а не на процесс: иначе процесс видит только свой GPU, и транспорт CUDA IPC не работает. Вариант \texttt{-{}-gpu-bind=closest} не используется, поскольку может выдать нескольким процессам один и тот же GPU.

\subsection{Уровень программы}

На уровне самой программы важен порядок вызовов. UCX читает свои параметры, в том числе сетевой адаптер (\texttt{UCX\_NET\_DEVICES}), при инициализации внутри вызова \texttt{MPI\_Init}, поэтому они должны быть заданы до него. Адаптер зависит от GPU процесса, а номер процесса внутри узла, по которому назначается GPU, до \texttt{MPI\_Init} доступен только из переменных окружения системы запуска (\texttt{SLURM\_LOCALID}).

\section{Способ измерения}

Задержка измеряется на стороне получателя таймером \texttt{MPI\_Wtime()}. Для упорядоченной пары GPU $(i,j)$ в $k$-м замере получатель фиксирует момент $t^{\mathrm{rdy}}_k$ готовности принять данные и момент $t^{\mathrm{end}}_k$ завершения приёма, а задержка оценивается как
\[
    t_{ij,k} = t^{\mathrm{end}}_k - t^{\mathrm{rdy}}_k.
\]
Завершение приёма имеет однозначный смысл: к этому моменту данные находятся в буфере получателя. Начальная отметка может предшествовать началу передачи, поэтому $t_{ij,k}$ является оценкой сверху для сквозной задержки. Зато обе отметки снимаются одним процессом, и согласование таймеров разных узлов не требуется.

Каждому GPU соответствует отдельный MPI-процесс. Open MPI используется в режиме CUDA-aware, в котором вызовы MPI принимают указатели на память GPU.

\textbf{Среды передачи.} В среде \texttt{auto} указатели на буферы в памяти GPU передаются непосредственно в вызовы MPI, а путь передачи выбирает UCX. После приёма выполняется синхронизация устройства, поскольку транспорт CUDA IPC может сообщить о завершении приёма раньше, чем копирование в память GPU фактически завершится. В среде \texttt{host} программа сама копирует данные в оперативную память и передаёт их средствами MPI, а для пар внутри одного узла использует общий буфер в разделяемой памяти узла.

\textbf{Режимы тестирования.} В режиме \texttt{one\_to\_one} пары GPU перебираются последовательно, и в каждый момент времени передача идёт только между одной парой. В режиме \texttt{all\_to\_all} передачи между всеми парами идут одновременно. Начальная отметка у получателя в этом режиме общая для всех источников, а конечная фиксируется отдельно для каждого источника.

\section{Экспериментальная среда и методика измерений}

Измерения выполнены на вычислительном комплексе ЦКП <<Информатика>>~[\Refn{ckp}] на узлах двух типов с GPU NVIDIA A100 40~ГБ. Узлы xFusion G5500 V7 содержат по четыре GPU, подключённых по шине PCIe, и по два процессора Intel Xeon Gold 6444Y. Узлы xFusion G8600 V7 содержат по восемь GPU в исполнении SXM4, объединённых коммутаторами NVSwitch, и по два процессора Intel Xeon Platinum 8462Y+. Узлы соединены сетью InfiniBand 100~Гбит/с.

В измерениях использовались по два GPU на двух узлах G5500 V7 и четыре GPU на одном узле G8600 V7. Таким образом, рассматриваются упорядоченные пары GPU трёх типов: 12 пар внутри узла G8600 V7 (NVLink), 4 пары внутри узлов G5500 V7 и 8 пар на разных узлах G5500 V7 (PCIe и InfiniBand). Теоретическая пропускная способность NVLink третьего поколения, установленного на A100, составляет 300~ГБ/с в одном направлении, а сети InfiniBand 100~Гбит/с~--- 12.5~ГБ/с.

Модуль сам выставляет переменные UCX, от которых зависит путь передачи: набор транспортов, разрешение GPUDirect RDMA и сетевой адаптер (например, \texttt{mlx5\_0:1}). Адаптер выбирается среди адаптеров InfiniBand по наиболее длинному общему префиксу пути в дереве устройств PCIe с путём к GPU, то есть за общим с ним коммутатором PCIe. Значения переменных записываются в заголовок каждого файла результатов. Переменная выставляется, только если она не задана в окружении, поэтому скрипт запуска перед каждым запуском сбрасывает значения, унаследованные из окружения пользователя.

Размер сообщения $S$ изменялся от 1~КБ до 16~МБ, и для каждой пары GPU и каждого размера выполнялась серия замеров. Результаты приводятся для режима \texttt{one\_to\_one}, в котором передачи разных пар не конкурируют между собой и задержка определяется путём передачи.

Первые передачи в каждой серии не учитываются: так из результата исключаются возможные разовые издержки на установление соединений транспортного слоя и регистрацию буферов. В замерах внутри серии могут наблюдаться периодические всплески задержки~[\Refn{volchaninov}], поэтому типичное значение для пары оценивается не средним, а медианой $\tilde t_{ij}$, устойчивой к выбросам~[\Refn{hoefler}]. Для каждого типа пар и размера сообщения приводится медиана по парам $\tilde t = \operatorname{med}_{(i,j)} \tilde t_{ij}$, а эффективная скорость передачи оценивается как $B = S/\tilde t$.

\section{Результаты}

\subsection{Влияние порога протокола UCX}

Поскольку протокол выбирается по размеру сообщения, в среде \texttt{auto} от размера зависит и путь передачи. Порог по умолчанию различается между группами узлов: на узле G8600 V7 он лежит между 8 и 16~КБ, на узлах G5500 V7~--- между 256 и 512~КБ. Ниже порога задержку ограничивает копирование через оперативную память, и при переходе через порог она резко падает (рис.~3).

\medskip
\noindent\begin{minipage}{\textwidth}
\includegraphics[width=\textwidth]{kibizov_421_thresh_ib}
\begin{center}
Рис.~3: Задержка для пар GPU на разных узлах G5500 V7 (PCIe + InfiniBand) в среде \texttt{auto} при трёх значениях \texttt{UCX\_RNDV\_THRESH} и в среде \texttt{host} при \texttt{UCX\_RNDV\_THRESH=0}, режим \texttt{one\_to\_one}
\end{center}
\end{minipage}
\medskip

В среде \texttt{host} порог на путь передачи не влияет: данные в любом случае проходят через оперативную память, и задержки при разных значениях порога различаются мало. Поэтому среда \texttt{host} далее сравнивается со средой \texttt{auto} при одном и том же пороге.

Если порог больше любого сообщения, все сообщения передаются протоколом eager, и для крупных сообщений передача оказывается в несколько раз медленнее, чем даже в среде \texttt{host}.

Поэтому основные измерения выполнены с \texttt{UCX\_RNDV\_THRESH=0}, при котором любое сообщение передаётся протоколом rendezvous и путь передачи не зависит от размера сообщения. Для небольших сообщений это почти не увеличивает задержку, а для сообщений среднего размера на узлах G5500 V7 заметно её снижает.

\subsection{Задержка и скорость передачи}

Медианные задержки при \texttt{UCX\_RNDV\_THRESH=0} для сообщений 1~КБ и 16~МБ приведены в таблице~2.

Для небольших сообщений задержка во всех случаях составляет десятки микросекунд, и заметно меньше она только в среде \texttt{auto} внутри узла G8600 V7. Для крупных сообщений задержку определяет скорость передачи по пути. В среде \texttt{auto} передача внутри узла G8600 V7 достигает около 75\% теоретической пропускной способности NVLink, а между узлами G5500 V7~--- около 70\% пропускной способности InfiniBand. Задержки отдельных пар одного типа различаются лишь на несколько процентов, то есть в режиме \texttt{one\_to\_one} неоднородность задержек определяется типом пары, а не отдельными парами.

\medskip
\noindent\begin{minipage}{\textwidth}
\begin{center}
\small
\renewcommand{\arraystretch}{1.15}
\begin{tabular}{@{}llrrr@{}}
\hline
 & & \multicolumn{2}{c}{$\tilde t$, мкс} & $B$, ГБ/с \\
Пары & Среда & 1~КБ & 16~МБ & 16~МБ \\
\hline
внутри G8600 & \texttt{auto} & 12.8 & 72.0 & 228 \\
 & \texttt{host} & 18.1 & 1258 & 13.0 \\
\hline
внутри G5500 & \texttt{auto} & 23.4 & 834 & 19.6 \\
 & \texttt{host} & 22.7 & 1253 & 13.1 \\
\hline
между G5500 & \texttt{auto} & 25.0 & 1839 & 8.9 \\
 & \texttt{host} & 23.1 & 2673 & 6.1 \\
\hline
\end{tabular}

\smallskip
{\normalsize Таблица~2: Медианная задержка и эффективная скорость передачи для пар внутри узла G8600 V7, внутри узла G5500 V7 и на разных узлах G5500 V7}
\end{center}
\end{minipage}
\medskip

Для пар внутри узла G5500 V7 скорость составляет лишь около 20~ГБ/с: по таблице выбора протоколов UCX (\texttt{UCX\_PROTO\_INFO=y}) в задании на двух узлах данные между GPU одного узла идут через сетевой адаптер. В задании на одном узле для тех же GPU используется транспорт \texttt{cuda\_ipc}, и 16~МБ передаются за 72~мкс, более чем в 10 раз быстрее. Как показали сообщения Open MPI, в задании на нескольких узлах процессы получают через систему запуска только сетевую часть адреса UCX друг друга, в том числе соседей по узлу, и другие транспорты UCX для них недоступны.

В среде \texttt{host} задержка для пар внутри узла одинакова на обеих группах узлов, то есть NVLink в этой среде не используется. Оценка по разности задержек для пар на разных узлах и внутри узла показывает, что сама передача по сети использует InfiniBand более чем на 90\%, то есть заметно лучше, чем GPUDirect RDMA.

\subsection{Выигрыш от передачи без промежуточного копирования}

На рис.~4 показано, во сколько раз передача в среде \texttt{auto} быстрее, чем в среде \texttt{host}. Внутри узла G8600 V7 среда \texttt{auto} выигрывает при любом размере сообщения, а для крупных сообщений~--- примерно в 17 раз.

На узлах G5500 V7 выигрыш появляется только для крупных сообщений: от 1~МБ между узлами и от 4~МБ внутри узла. Он не превышает 1.5 раза. Для пар внутри узла это объясняется тем, что данные идут через сетевой адаптер (раздел~5.2): у такого пути большая постоянная задержка, и для сообщений среднего размера среда \texttt{auto} оказывается даже медленнее среды \texttt{host}. Таким образом, передача без промежуточного копирования выгодна не всегда.

\medskip
\noindent\begin{minipage}{\textwidth}
\includegraphics[width=\textwidth]{kibizov_421_speedup}
\begin{center}
Рис.~4: Во сколько раз передача в среде \texttt{auto} быстрее, чем в среде \texttt{host} (отношение задержек \texttt{host}/\texttt{auto}), при \texttt{UCX\_RNDV\_THRESH=0}, режим \texttt{one\_to\_one}
\end{center}
\end{minipage}
\medskip

\section*{Заключение}

Разработан модуль, измеряющий задержки передачи данных между всеми упорядоченными парами GPU кластера. Путь передачи определяется не только аппаратурой, но и настройками Open MPI, UCX и Slurm, причём при неподходящих настройках UCX не сообщает об ошибке, а передаёт данные через оперативную память. Поэтому параметры UCX, включая порог протокола rendezvous, следует задавать явно и сохранять вместе с результатами, сетевой адаптер выбирать по близости к GPU, а ядра процессора выделять на сокете GPU.

Измерения показали, что выигрыш от передачи без промежуточного копирования зависит от расположения GPU и от пути, выбранного UCX: внутри узла G8600 V7 он наблюдается при любом размере сообщения, а между узлами G5500 V7~--- только для крупных сообщений. В задании на нескольких узлах UCX может передавать данные между GPU одного узла через сетевой адаптер, что замедляет передачу более чем в 10 раз. Поэтому выбранный UCX путь следует проверять, например с помощью \texttt{UCX\_PROTO\_INFO}.

В дальнейшем планируется исследовать режим \texttt{all\_to\_all} и выделить группы пар GPU со сходными задержками аналогично~[\Refn{salnikov2017}] на большем числе узлов.

\References
\begin{enumerate}
\item
    \Label{salnikov2017}
    Сальников~А.~Н., Бегаев~А.~А., Майсурадзе~А.~И.
    \emph{Кластеризация задержек между узлами суперкомпьютера и визуализация полученных результатов} // Современные проблемы математического моделирования, обработки изображений и параллельных вычислений 2017: труды Международной научной конференции (пос. Дивноморское, 4--11 сентября 2017~г.). Ростов-на-Дону: ДГТУ-Принт, 2017. Т.~1. С.~206--215.

\item
    \Label{begaev2020}
    Бегаев~А.~А., Сальников~А.~Н.
    \emph{Метод измерения задержек при передаче данных между графическими процессорами, находящимися на разных узлах вычислительного кластера} // Вестник Московского университета. Серия~15: Вычислительная математика и кибернетика. 2020. Вып.~1. С.~3--13.

\item
    \Label{clustbench}
    \emph{Clustbench / network-tests2, ветка gpu\_tests.}
    \url{https://github.com/clustbench/network-tests2/tree/gpu_tests} (дата обращения: 28.09.2026).

\item
    \Label{cudap2p}
    NVIDIA. \emph{CUDA C++ Programming Guide: Peer Device Memory Access.}
    \url{https://docs.nvidia.com/cuda/cuda-c-programming-guide/} (дата обращения: 20.09.2026).

\item
    \Label{gpudirect}
    NVIDIA. \emph{GPUDirect RDMA.}
    \url{https://docs.nvidia.com/cuda/gpudirect-rdma/} (дата обращения: 20.09.2026).

\item
    \Label{openmpi}
    \emph{Open MPI 5.0.x Documentation: CUDA.}
    \url{https://docs.open-mpi.org/en/v5.0.x/tuning-apps/networking/cuda.html} (дата обращения: 20.09.2026).

\item
    \Label{ucx}
    \emph{Unified Communication X.}
    \url{https://openucx.org/} (дата обращения: 20.09.2026).

\item
    \Label{slurm}
    \emph{Slurm Workload Manager: Generic Resource (GRES) Scheduling.}
    \url{https://slurm.schedmd.com/gres.html} (дата обращения: 20.09.2026).

\item
    \Label{ckp}
    \emph{Центр коллективного пользования <<Информатика>> ФИЦ ИУ РАН: Вычислительная инфраструктура.}
    \url{https://docs.hpc.frccsc.ru/hardware/} (дата обращения: 20.09.2026).

\item
    \Label{volchaninov}
    Волчанинов~А.~П., Сальников~А.~Н.
    \emph{Исследование структуры серии измерений задержки при нагрузочном тестировании коммутационной среды вычислительного кластера} // Суперкомпьютерные дни в России: труды международной конференции (Москва, 25--26 сентября 2023~г.). М.: МАКС Пресс, 2023. С.~42--52.

\item
    \Label{hoefler}
    Hoefler~T., Belli~R.
    \emph{Scientific Benchmarking of Parallel Computing Systems: Twelve Ways to Tell the Masses when Reporting Performance Results} // Proceedings of SC'15. 2015. Art.~73. P.~1--12. DOI: 10.1145/2807591.2807644.

\end{enumerate}
